% Separate post-CP18 candidate. This does not assert a Petri formula for D_t.
\subsection{A norm pencil on the determinant curve}

\begin{lemma}[The determinant cokernel is residual to a squared pencil]
\label{lem:DE-A3-positive-pencil-candidate}
\status{Candidate; complete local proof, independent audit pending.}
Keep the algebraically closed characteristic-two hypotheses and notation of
Lemma~\ref{lem:DE-A3-positive-determinant-curve}, and take a general pair
with disjoint zero schemes. On $D=D_s$ the images of $s_0$ and $s_1$
generate a line subbundle $R_0\subset E_0|_D$, after viewing the latter
section as a map $P^{-1}|_D\to E_0|_D$. Then
\[
 R_0\simeq\mathcal O_D(Z_0),\qquad
 R_0P\simeq\mathcal O_D(Z_1),\qquad
 P|_D\simeq\mathcal O_D(Z_1-Z_0).
\]
In particular $Z_0,Z_1$ are effective Cartier divisors on $D$, including
when $D$ is nonreduced. Put $A=R_0^{\otimes2}$. The squares of the two
sections, using the fixed trivialization of $P^2$, give a base-point-free
pencil $\langle a_0,a_1\rangle\subset H^0(D,A)$ with divisors $2Z_0,2Z_1$.
The line bundle in (DE-TB1) is
\[
                         R_s\simeq\omega_D\otimes A^{-1}.
                                                        \tag{DE-TP1}
\]
Here $\deg A=b$ and $p_a(D)=b+1$.
\end{lemma}
\begin{proof}
At every point of $D$, at least one of $s_0,s_1$ is nonzero, and their
determinant vanishes. In a frame with the nonvanishing section first,
their matrix has a unit entry and determinant zero. Its image on $D$
is therefore a line subbundle; this constructs $R_0$ even over the
nonreduced local rings of $D$. Near $Z_0$, use $s_1$ to trivialize
the common line. The two equations of $Z_0$ on the smooth surface can
then be written $(f,w)$, where $w$ cuts out $D$. They form a regular
sequence. Consequently $f$ is a nonzerodivisor on $D$ and cuts out
$Z_0$ there. The same argument treats $Z_1$. The asserted line-bundle
identifications follow from their sections. Their squares have no common
zero, since $Z_0\cap Z_1=\varnothing$.

Tensor (DE-TB1) by $\mathcal O_D$. Since $D$ is Cartier, its initial
Tor term is $R_s\otimes\mathcal O_D(-D)$, and hence
\[
 R_s(-D)=\ker\bigl(\mathcal O_D^2\longrightarrow Q_s|_D\bigr).
\]
The earlier local chart has columns $(1,0)$ and $(u f^2,w)$,
where $e^2\mapsto u$ for the chosen frame $e$ of $P$.
On $D$ their kernel is the kernel of the pencil row $(1,u f^2)$.
Globally this row is $(a_0,a_1):\mathcal O_D^2\to A$, whose kernel
is $A^{-1}$. Thus $R_s(-D)=A^{-1}$. Adjunction on the abelian surface
identifies $\mathcal O_D(D)$ with $\omega_D$, proving (DE-TP1).
The degree statements follow from $\operatorname{length}Z_i=b/2$ and
the previously computed arithmetic genus. No reducedness or smoothness
of $D$ was used.
\end{proof}

\begin{lemma}[Norm factorization and the actual ordinary differential]
\label{lem:DE-A3-positive-norm-pencil-candidate}
\status{Candidate; complete local proof, independent audit pending.}
With the same general pair, the norm of the $\alpha_2$-quotient
$\nu:C\to T$ gives a section
\[
                   \sigma\in H^0(D,R_sP)
\]
which is a nonzerodivisor on $D$. If $S_i\in H^0(T,L)$ is the
Frobenius-twisted section corresponding to $s_i$, then
\[
                         S_i|_D=\sigma a_i.             \tag{DE-TP2}
\]
There is a canonical isomorphism
\[
 \frac{H^0(T,L)}{\langle S_0,S_1\rangle}
 \simeq
 \frac{H^0(D,\omega_DP)}{\sigma\langle a_0,a_1\rangle}.
                                                        \tag{DE-TP3}
\]
The fixed trivialization of $P^2$ also defines a nonzero invariant
differential $\eta_P\in H^0(T,\Omega_T^1)$ with
$C(\eta_P)=\eta_P$. Wherever $a_0a_1$ is invertible on $D$,
\[
                      d\log(a_1/a_0)=\eta_P|_D.         \tag{DE-TP4}
\]
This identity does not assert that its right side is nonzero on $D$.
\end{lemma}
\begin{proof}
Locally the finite flat $\alpha_2$-algebra is
$\mathcal O_T[z]/(z^2-h)$. In a trivialization of $J$, its norm is
\[
                         n(a+bz)=a^2+h b^2.
\]
Determinants of the transition matrices identify its target with
$\det\nu_*J=L$; the determinant of $\nu_*\mathcal O_C$ is trivial,
as the torsor algebra is an extension of $\mathcal O_T$ by itself.
Thus this construction gives the intrinsic quadratic norm
$n:E_0\to L$. For $E_1=E_0P$, the same norm is multiplied by the
local unit $u$ specifying $P^2=\mathcal O_T$. This is precisely the
descent of the square of the original section through relative
Frobenius, so $n(s_i)=S_i$.

The restriction of the norm to the line $R_0$ is a homomorphism
$A=R_0^2\to L|_D$. It is therefore a section of
$L|_D\otimes A^{-1}=R_sP$, using $\omega_D=(LP)|_D$.
This is $\sigma$, and homogeneity gives (DE-TP2). For a general
pair, the divisors of the sections of $J$ and $JQ$ on $C$ have no
common component: both line bundles are globally generated, and
one may choose the second section to avoid the finitely many
components of the first. Their norm divisors on $T$ likewise have
no common component, since $\nu$ is a homeomorphism on underlying
spaces. If $\sigma$ belonged to a minimal prime of $D$, (DE-TP2)
would put both norm sections in the corresponding prime. This is
impossible. The Cartier divisor $D$ is Cohen--Macaulay and has no
embedded associated primes, so $\sigma$ is a nonzerodivisor.

The restriction sequence is
\[
                 0\longrightarrow P\longrightarrow L
                    \longrightarrow L|_D\longrightarrow0.
\]
Acyclicity of $P$ identifies $H^0(T,L)$ with $H^0(D,L|_D)$.
As $L|_D=\omega_DP$, equation (DE-TP2) gives (DE-TP3).

Finally, write $e^2\mapsto u$ in any local frame of $P$.
On overlaps $u$ changes by the square of a transition unit,
so $d\log u$ glues to a global differential $\eta_P$.
The usual Cartier identity gives $C(d\log u)=d\log u$.
If this differential were zero, the Cartier kernel calculation
on the smooth perfect-base surface would make every $u$ locally
a square. The uniquely determined square roots then normalize
all transition functions to one, trivializing $P$, a contradiction.
Every global differential on an abelian variety is invariant.
In the determinant chart the pencil ratio is exactly $u f^2$,
so its logarithmic differential is $d\log u$. This proves
(DE-TP4), with no assumption that $u$ itself has a local square root.
\end{proof}

\begin{remark}[What this model has not established]
\label{rem:DE-A3-positive-norm-pencil-scope-candidate}
\status{Candidate scope statement; the quadratic variation remains unresolved.}
The two lemmas identify the quotient domain of the variation with an
explicit twisted canonical space on a Cartier curve and retain both
the norm factor and the ordinary logarithmic differential. They do
not identify the matrix-valued variation with a Petri or residue map.
In particular no conclusion $h^0(A)=2$, no smoothness of $D$, no
generic separability of its pencil, and no lower bound of three on
the variation-image dimension is asserted. Those conclusions would
require further arguments. The common-kernel/common-image exclusion
remains the open step.
\end{remark}
